\documentclass[titlepage,a4paper]{jsarticle}
\usepackage{../../sty/import}% 各種パッケージインポート
\usepackage{../../sty/title}% タイトルページの変更

%% タイトルページの変数
% レポートタイトル
% ダディダディ
\title{理論生命科学最終レポート}
% 提出日
\expdate{\today}
% 科目名
\subject{理論生命科学}
% 分野
\class{情報経営システム工学分野}
% 学年
\grade{M1}
% 学籍番号
\mynumber{24336488}
% 記述者
\author{本間三暉}

\begin{document}
% titleページ作成
\maketitle
\section{Find a self-organization system around you.Whichever everyday example or professional example.
Describe how the example you chose is a self-organizing system.}
% 身の回りにある自己組織化システムを1つ見つけなさい。
% ※ 日常生活の例でも，専門分野の例でもよい。
% 選んだ例がどのように自己組織化システムであるかを説明しなさい。
I chose the spread of memes on social networking services as an example of a self-organizing system.
A meme is an image, phrase, video, or style of expression that is repeatedly shared and modified by many users.

Each user acts independently and does not control the entire trend. 
Users simply react to posts that they find interesting. 
For example, they may like a post, share it with other users, imitate its format, or create a modified version. 
Other users then see these reactions and decide whether to participate in the same way.
Through these repeated local interactions, a meme can rapidly spread across the whole social network.

As more users participate, common phrases, images, and patterns of expression naturally emerge.
 No single person designs the complete structure of the trend or gives instructions to all users. Instead, a large-scale trend is created from the independent decisions and interactions of many individuals.

Recommendation algorithms and influential users can affect how quickly a meme spreads. However, they cannot completely determine how every user reacts or how the meme changes. Therefore, the spread of memes on social networking services can be regarded as a self-organizing system in which individual behavior produces an organized collective pattern.


\section{Is this good example of self-organizing system or not? Why do you think so?}
% 質問：これは自己組織化システムの良い例であるか，そうでないか．なぜそのように考えるのか説明しなさい．
I do not think this is a good example of a self-organizing system. In a self-organizing system, each individual acts autonomously and interacts with nearby individuals. As a result of these local interactions, an organized pattern emerges without a central leader controlling the whole group.

However, the collective behavior shown in the video appears to be planned and learned through training. 
Each individual follows predetermined movements, positions, or instructions, rather than deciding how to act based only on local interactions. 
The overall pattern is therefore created by top-down control.

Although the group moves in an organized manner, organized collective behavior is not always self-organization. 
Since the behavior is intentionally designed and controlled in advance, I do not consider it a good example of a self-organizing system.

\section{Explain the methods to induce the Rubber Hand Illusion.}
% ラバーハンド錯覚を引き起こす方法について説明しなさい．
The rubber hand illusion can be induced by using a rubber hand, a screen, and two brushes. 
First, the participant places one hand on a table, and the real hand is hidden from view behind a screen. 
A rubber hand is then placed in front of the participant in a natural position close to the hidden real hand. The participant is asked to look at the rubber hand.

Next, the experimenter strokes the hidden real hand and the visible rubber hand at the same time, in the same location and direction. 
This synchronous stimulation is repeated for several minutes. 
Because the participant sees the rubber hand being touched while feeling the same touch on the real hand, the brain integrates the visual and tactile information.

As a result, the participant may begin to feel that the rubber hand is part of their own body. 
The illusion becomes weaker or may not occur when the two hands are touched at different times or when the rubber hand is placed in an unnatural position.
\section{Define the word "Phantom pain"}
% ファントムペインの定義を説明しなさい．(幻肢痛)
Phantom pain is pain that a person feels in a body part that has been amputated or is no longer physically present. 
It most commonly occurs after the loss of a limb.

\section{かだい５}
% 参考文献
\begin{thebibliography}{99}
\bibitem{授業スライド}理論生命科学 授業スライド
\end{thebibliography}

\end{document}